\chapter{Chapter06 — Interpolation}
%%% generated by the blueprint pipeline from FABL.Chapter06.F₂Polynomials.Interpolation
\section{Overview}

Declarations of the Lean module \texttt{FABL.Chapter06.F₂Polynomials.Interpolation}.

\section{Declarations}

\begin{definition}[Binary cube point indicator]
\lean{FABL.f₂PointIndicator}
\leanok
\uses{FABL.F₂Cube, FABL.𝔽₂}
\statementsource{Analysis-of-Boolean-Functions-by-Ryan-ODonnell}{
  pdf-ffbe9c7b3306-p021-b002,
  pdf-ffbe9c7b3306-p148-b005
}
For $n\in\mathbb{N}$, let $a,x\colon\{0,1\}^n\to\mathbb{F}_2$, where $\mathbb{F}_2$ is
the field with two elements. Define the indicator of $a$ evaluated at $x$ by
$\displaystyle \prod_{i=0}^{n-1}\begin{cases}x_i,&a_i=1,\1-x_i,&a_i\ne 1.\end{cases}$
\end{definition}

\begin{theorem}[F₂point indicator eq ite]
\lean{FABL.f₂PointIndicator_eq_ite}
\leanok
\uses{FABL.F₂Cube, FABL.f₂PointIndicator}
\statementsource{Analysis-of-Boolean-Functions-by-Ryan-ODonnell}{pdf-ffbe9c7b3306-p148-b005}
\difficulty{5}
% dataset-pipeline: stub -- prose pending build_informal_statements.py
Statement pending regeneration.
\end{theorem}

\begin{theorem}[Interpolation by point indicators]
\lean{FABL.f₂Interpolation}
\leanok
\uses{FABL.F₂BooleanFunction, FABL.F₂Cube, FABL.f₂PointIndicator, FABL.f₂PointIndicator_eq_ite}
\difficulty{3}
\proofsource{Analysis-of-Boolean-Functions-by-Ryan-ODonnell}{
  pdf-ffbe9c7b3306-p021-b002,
  pdf-ffbe9c7b3306-p021-b004,
  pdf-ffbe9c7b3306-p149-b001,
  pdf-ffbe9c7b3306-p149-b002
}
\proofstep
  {FABL.𝔽₂}
  {context}
  {Analysis-of-Boolean-Functions-by-Ryan-ODonnell}
  {pdf-ffbe9c7b3306-p015-b008,
    pdf-ffbe9c7b3306-p148-b003}
\proofstep
  {FABL.F₂Cube}
  {context}
  {Analysis-of-Boolean-Functions-by-Ryan-ODonnell}
  {pdf-ffbe9c7b3306-p028-b004,
    pdf-ffbe9c7b3306-p148-b003}
\proofstep
  {FABL.F₂BooleanFunction}
  {context}
  {Analysis-of-Boolean-Functions-by-Ryan-ODonnell}
  {pdf-ffbe9c7b3306-p149-b002,
    pdf-ffbe9c7b3306-p149-b004}
\proofstep
  {FABL.f₂PointIndicator}
  {direct}
  {Analysis-of-Boolean-Functions-by-Ryan-ODonnell}
  {pdf-ffbe9c7b3306-p148-b005}
\proofstep
  {FABL.f₂PointIndicator_eq_ite}
  {context}
  {Analysis-of-Boolean-Functions-by-Ryan-ODonnell}
  {pdf-ffbe9c7b3306-p148-b005}
Let $n$ be a natural number, let $f:\mathbb F_2^n\to\mathbb F_2$ be any function, and
let $x\in\mathbb F_2^n$. The sum over all $a\in\mathbb F_2^n$ of $f(a)$ times the point
indicator of $a$ evaluated at $x$ equals $f(x)$.
\end{theorem}

\begin{definition}[Constant-input indicator]
\lean{FABL.f₂EqualityFunction}
\leanok
\uses{FABL.F₂Cube, FABL.𝔽₂}
\statementsource{Analysis-of-Boolean-Functions-by-Ryan-ODonnell}{pdf-ffbe9c7b3306-p148-b005}
Define a function on the $n$-dimensional binary cube, valued in the two-element field,
to be $1$ when all coordinates of its input are equal and $0$ otherwise.
\end{definition}

\begin{theorem}[Constant binary vectors]
\lean{FABL.forall_coordinate_eq_iff_eq_zero_or_eq_one}
\label{FABL.forall_coordinate_eq_iff_eq_zero_or_eq_one}
\leanok
\uses{FABL.F₂Cube}
\difficulty{4}
Let $n>0$ and let $x$ be a vector of length $n$ over the field with two elements. All
coordinates of $x$ are equal if and only if $x$ is the zero vector or the vector whose
coordinates are all one.
\end{theorem}

\begin{theorem}[Equality function as a sum of point indicators]
\lean{FABL.f₂EqualityFunction_eq_pointIndicators}
\leanok
\uses{FABL.F₂Cube, FABL.forall_coordinate_eq_iff_eq_zero_or_eq_one, FABL.f₂EqualityFunction, FABL.f₂PointIndicator, FABL.f₂PointIndicator_eq_ite}
\difficulty{5}
Let $n$ be a positive natural number and let $x\in\mathbb F_2^n$. The function that is
$1$ exactly when all coordinates of $x$ are equal, and $0$ otherwise, equals the sum
over $\mathbb F_2$ of the point indicators of the all-zero and all-one vectors,
evaluated at $x$.
\end{theorem}

\begin{theorem}[Three-bit parity polynomial identity]
\lean{FABL.threeBitParity_integer_interpolation}
\proofsource{Analysis-of-Boolean-Functions-by-Ryan-ODonnell}{
  pdf-ffbe9c7b3306-p035-b010,
  pdf-ffbe9c7b3306-p149-b001
}
\label{FABL.threeBitParity_integer_interpolation}
\leanok
\difficulty{2}
For all integers $x_1,x_2,x_3$,
\[(1-x_1)(1-x_2)x_3+(1-x_1)x_2(1-x_3)+x_1(1-x_2)(1-x_3)+x_1x_2x_3=x_1+x_2+x_3-2(x_1x_2+x_1x_3+x_2x_3)+4x_1x_2x_3.\]
\end{theorem}

\begin{theorem}[Three-bit parity interpolation over the binary field]
\lean{FABL.threeBitParity_f₂_interpolation}
\leanok
\uses{FABL.𝔽₂}
\proofsource{Analysis-of-Boolean-Functions-by-Ryan-ODonnell}{pdf-ffbe9c7b3306-p149-b001}
\proofstep
  {FABL.𝔽₂}
  {context}
  {Analysis-of-Boolean-Functions-by-Ryan-ODonnell}
  {pdf-ffbe9c7b3306-p015-b008}
\difficulty{3}
For $x_1,x_2,x_3\in\mathbb F_2$, the sum of the four products that indicate the
odd-parity inputs equals $x_1+x_2+x_3$ in $\mathbb F_2$:
\[(1-x_1)(1-x_2)x_3+(1-x_1)x_2(1-x_3)+x_1(1-x_2)(1-x_3)+x_1x_2x_3=x_1+x_2+x_3.\]
\end{theorem}
